% !TeX program = pdflatex
%
% market-positioning-paper.tex
%
% Build:  pdflatex market-positioning-paper.tex   (run 3x: ToC, then cross-refs)
%

\documentclass[11pt,a4paper]{article}

\usepackage[T1]{fontenc}
\usepackage[utf8]{inputenc}
\usepackage[margin=2.4cm]{geometry}
\usepackage{amsmath}
\usepackage{amssymb}
\usepackage{booktabs}
\usepackage{longtable}
\usepackage{array}
\usepackage{xcolor}
\usepackage{enumitem}
\usepackage{fancyhdr}
\usepackage{ragged2e}
\usepackage[hidelinks]{hyperref}

\definecolor{warn}{RGB}{160,60,0}
\definecolor{ok}{RGB}{20,110,60}
\definecolor{mid}{RGB}{120,95,0}
\newcommand{\bad}[1]{\textcolor{warn}{#1}}
\newcommand{\good}[1]{\textcolor{ok}{#1}}
\newcommand{\mixed}[1]{\textcolor{mid}{#1}}
\newcommand{\code}[1]{\texttt{\small #1}}
\newcommand{\name}[1]{\textsc{#1}}

\newcolumntype{L}[1]{>{\RaggedRight\arraybackslash}p{#1}}

\pagestyle{fancy}
\fancyhf{}
\lhead{\small xT Market Positioning}
\rhead{\small\thepage}
\renewcommand{\headrulewidth}{0.4pt}

\title{\textbf{Positioning a Pharmacokinetic Instrument in a\\
Hostile and Under-Served Market:}\\[0.3em]
\large A Market, Cultural, and Nomenclature Study of xT,\\
with Screened Rebrand Candidates}
\author{Prepared for Starla Cortes}
\date{August 2026}

\begin{document}
\maketitle

\begin{abstract}
\noindent
xT is a mobile application that predicts serum hormone concentration
over time from a user's logged dose history for gender-affirming
hormone replacement therapy (HRT), built on a superposed Bateman
biexponential model with per-user calibration. Its engineering is
documented to a standard unusual for a consumer application. Its name
is a development codename. This paper asks what market the product is
actually entering, what that market's cultural and political conditions
demand of a brand, and what the application should be called.

We find three things that revise the project's own prior positioning
statements. First, the trans-specific segment of this category is not
competitive but \emph{vacant}: its most-cited entrants---Transcapsule,
Solace/Euphoria, MyTransHealth---are abandoned, with last updates in
2022, 2022, and 2019 respectively, while the actively maintained,
well-rated products in the adjacent space are coded to
cisgender male optimization (Regimen, 338 ratings, 4.8, updated within
the week) or GLP-1 weight loss (DoneDose). The competitive threat is
not feature parity; it is that this category has taught its users that
tools built for them get abandoned. Second, and correcting an explicit
claim in the project's own business plan, the pharmacokinetic model is
\emph{not} a unique differentiator: Regimen ships a public one-compartment
Bateman superposition visualizer with cited sources and stated
limitations. xT's defensible asset is not that it models
pharmacokinetics but the depth and honesty of its disclosure---a
published derivation, an out-of-sample validation dossier, and a
revision history that names its own defects. Third, the brand brief the
founder has set---clinically credible \emph{and} discreet---is not a
stylistic preference but a safety requirement, and it is the specific
intersection at which every prior entrant in this category has failed.

We evaluate rebrand candidates against a four-part discretion test
(lock screen, financial record, over-the-shoulder, and disclosure under
compulsion) and a clearance screen covering App Store collisions and
United States trademark records in classes 9, 42, and 44. We report
seven screened candidates and eleven rejections with reasons, and we
benchmark the incumbent codename against them. We recommend
\name{Isocline}, with \name{Ephemeris} as the principal alternate and a
defensible case for retaining \name{xT} stated in full rather than
strawmanned.

Clearance findings in this paper are indicative, not legal opinions.
Where a source is a competitor's own marketing, we say so.

\textbf{Keywords:} market positioning, brand nomenclature, trademark
screening, gender-affirming care, digital health, app store
optimization, health data privacy, trans health technology
\end{abstract}

\tableofcontents
\clearpage

% ======================================================================
\section{Introduction}
\label{sec:intro}

\subsection{The problem this paper addresses}

xT began as a directory name. It is short, it gestures at testosterone,
it is at least ungendered in the sense that it does not commit the
product to one hormone---but nobody chose it, and no part of it was
tested against the conditions the product will actually ship into. The
question posed for this study was whether xT should be renamed, and if
so, to what.

That question cannot be answered by generating names. A name is a
solution to a positioning problem, and the positioning problem here is
unusual in three ways that a generic naming exercise would miss:

\begin{enumerate}[leftmargin=1.5em,itemsep=0.3em]
\item The product's technical asset is a pharmacokinetic model of
      unusual rigor for a consumer application, which means the brand
      has to be able to carry a scientific claim without over-claiming
      a regulatory one.
\item The user population is, in a material and documented sense, under
      surveillance. A name is a string that appears on a lock screen, a
      bank statement, and a subpoenaed app-store purchase record. In
      this category that is a safety property, not a marketing one.
\item The category has a mortality problem. The applications this
      product would be compared against are, disproportionately,
      dead. That fact shapes what a new entrant's brand has to
      communicate more than any feature comparison does.
\end{enumerate}

\subsection{Scope}

The founder scoped this study to three of four candidate research
areas: the competitive application landscape
(\S\ref{sec:landscape}), the social, cultural, and political
environment (\S\ref{sec:sociopolitical}), and brand and category
positioning (\S\ref{sec:positioning}). Medical and clinical positioning
---provider adoption, WPATH and Endocrine Society context, HIPAA
posture, and the software-as-a-medical-device question---was explicitly
deprioritized and is not covered here.

One exception is taken to that exclusion, and it is flagged rather than
smuggled. In March 2026 Apple began requiring apps in the Health \&
Fitness and Medical categories to declare a regulated-medical-device
status, with an early-2027 deadline after which undeclared apps cannot
ship updates~\cite{appledevice}. This is a regulatory fact, but it
constrains what the brand is permitted to \emph{claim}, and therefore
falls inside the positioning scope rather than outside it. It is
treated in \S\ref{sec:deviceclaim} and nowhere else.

\subsection{The brand brief}

The founder specified two optimization targets, selected together:

\begin{description}[leftmargin=1.6em,style=nextline,itemsep=0.35em]
\item[Clinically credible.] The name should read as a serious health
      instrument---provider-shareable, grant-fundable, and able to sit
      on the title page of a research paper without embarrassment.
\item[Discreet and ungendered.] The name should disclose nothing on a
      lock screen or a financial record.
\end{description}

These are not independent axes, and the fact that they were selected
together is the most consequential input to this study. Section
\ref{sec:brief} shows that the intersection of the two is a narrow
region that no existing product in this category occupies, that the
occupied regions are occupied for legible reasons, and that most of the
obvious naming instincts---including several this study generated and
then discarded---fail one axis in the course of satisfying the other.

\subsection{Contributions}

\begin{enumerate}[leftmargin=1.5em,itemsep=0.25em]
\item A competitive census of the HRT-tracking category built from
      primary App Store records rather than from competitor-authored
      comparison articles (\S\ref{sec:census}), including
      last-updated dates, rating counts, and pricing.
\item An abandonment analysis showing that the trans-specific segment
      is vacant rather than contested, and an argument that this
      changes what the brand must do (\S\ref{sec:graveyard}).
\item A correction to the project's own business plan: the claim that
      no competitor models pharmacokinetics is false
      (\S\ref{sec:notamoat}).
\item A four-part discretion test operationalizing the ``discreet''
      half of the brief (\S\ref{sec:discretion}).
\item Seven screened rebrand candidates and eleven documented
      rejections, with App Store and trademark screens
      (\S\ref{sec:candidates}, \S\ref{sec:rejected}).
\item A stated case for retaining the incumbent codename, benchmarked
      against the candidates rather than dismissed
      (\S\ref{sec:incumbent}).
\end{enumerate}

% ======================================================================
\section{Method and source quality}
\label{sec:method}

\subsection{Why this section exists}

The pharmacokinetic paper accompanying this project grades its own
constants---primary-sourced, secondary-sourced, or uncalibrated---and
declines to present a secondary source as though it were a primary one.
Market research is more prone to the same failure, not less, because
its most available sources are written by interested parties. This
paper adopts the same hierarchy.

\subsection{Source tiers}

\begin{longtable}{@{}L{2.6cm}L{5.4cm}L{7.2cm}@{}}
\toprule
\textbf{Tier} & \textbf{Definition} & \textbf{Used in this paper for} \\
\midrule
\endhead
Primary &
Record produced by the party with direct knowledge, published for a
purpose other than persuading the reader of a market conclusion. &
App Store listing pages (pricing, rating counts, last-updated dates,
seller names, declared data practices); USPTO records via Justia;
Apple developer documentation and policy announcements. \\[0.4em]

Secondary &
Independent third-party reporting or research with no commercial
interest in the category. &
Peer-reviewed and press reporting on DIY HRT prevalence; the Movement
Advancement Project policy maps; POGO's surveillance investigation;
appellate reporting. \\[0.4em]

\bad{Interested} &
Written by a market participant about its own market. &
The ``best HRT tracker apps'' comparison articles published by Regimen
and DoneDose. Used only for feature enumeration, never for ranking or
evaluation. \\[0.4em]

\bad{Unverified} &
Asserted in a project document without a traceable source, or resistant
to verification with the tools available. &
Domain availability; the completeness of any trademark screen; several
claims in the project's own business plan, one of which is corrected in
\S\ref{sec:notamoat}. \\
\bottomrule
\end{longtable}

\subsection{A note on the comparison articles}

Two of the most prominent search results for ``best HRT tracker apps''
are \code{helloregimen.com/blog/\allowbreak best-hrt-tracker-apps-2026}
and \code{donedose.com/guides/\allowbreak best-hrt-tracker-app}. Both are
content
marketing published by applications that appear in their own
rankings~\cite{regimenblog,donedoseblog}. Regimen's article places
Regimen first. DoneDose's article places DoneDose first. Neither
discloses the conflict in the body text.

This is not a scandal; it is the normal structure of search-optimized
category content, and both articles are factually useful for
enumerating who exists. But it has a specific consequence for this
project: the competitive picture in the project's existing business
plan is built substantially on these two articles, and therefore
inherits their framing. Where this paper disagrees with that plan, the
disagreement usually traces to having gone to the App Store listing
instead.

\subsection{What this paper cannot establish}
\label{sec:cannot}

Three limits are stated here rather than buried in
\S\ref{sec:limitations}, because they bear on how the rebrand section
should be read.

\textbf{Trademark screens here are not clearance opinions.} The screens
in \S\ref{sec:candidates} were run against Justia's index of United
States federal trademark records, filtered to International Classes 9
(software), 42 (technology services), and 44 (medical services). They
do not cover state registrations, unregistered common-law marks,
foreign registrations, or the likelihood-of-confusion analysis that
determines whether a mark is actually available. A mark reported here
as ``clear'' means only that this screen surfaced no conflict in those
classes. Retaining a trademark attorney before filing or before
printing anything is a separate and necessary step.

\textbf{Domain availability was not established.} The tooling available
for this study could not query a registrar. Domain notes in
\S\ref{sec:candidates} are inferences from whether a site appears to
exist, and should be re-checked directly---a task of a few minutes per
name.

\textbf{App Store name availability was not established.} Whether a
name can be reserved is determined in App Store Connect, not by
searching the storefront. The App Store screens here establish only
whether a user searching that word would encounter competing results.

% ======================================================================
\section{The competitive landscape}
\label{sec:landscape}

\subsection{Three overlapping markets, not one}
\label{sec:threemarkets}

``HRT tracking'' names three populations that share a mechanism and
share almost nothing else. Conflating them is the single most common
error in this category's own marketing, and it is why the comparison
articles read as incoherent---they rank a menopause education platform
against an injection logger against a peptide tracker.

\begin{longtable}{@{}L{3.2cm}L{5.0cm}L{3.1cm}L{3.6cm}@{}}
\toprule
\textbf{Market} & \textbf{Who} & \textbf{Emotional register} &
\textbf{Representative apps} \\
\midrule
\endhead
Menopausal HRT &
Cisgender women, typically 45--60, usually clinician-managed, often
national-health-service adjacent. &
Symptom relief; being taken seriously. &
Balance, Health \& Her, Caria, mySysters, Hormona, HRTMe \\[0.4em]

Optimization / TRT / peptides &
Predominantly cisgender men, self-directed, protocol-tinkering,
overlapping with GLP-1 and bodybuilding communities. &
Performance; control; data maximalism. &
Regimen, DoneDose, TRT Tracker \\[0.4em]

Gender-affirming HRT &
Trans and gender-diverse adults, frequently self-advocating, sometimes
unsupervised, materially at legal and social risk. &
Safety; self-determination; not being surveilled. &
Transcapsule, Solace, Hormone Helper, Trans Memo \\
\bottomrule
\end{longtable}

The three markets differ in what a product must earn. The menopause
market rewards authority---Balance leads with its founder's clinical
credentials and its alignment with national guidance~\cite{regimenblog}.
The optimization market rewards granularity and rewards it immediately;
Regimen's pitch is 200-plus compounds and 50-plus daily
markers~\cite{regimenlisting}. The gender-affirming market rewards
neither of those first. It rewards evidence that the builder is not
going to disappear, extract, or expose them---and it has extensive
evidence that builders do all three.

\subsection{Census of the actively relevant field}
\label{sec:census}

The following is built from App Store listing pages, which are primary
records, and not from the comparison articles. Figures were retrieved
in August 2026 and will drift.

\begin{longtable}{@{}L{3.1cm}L{2.5cm}L{1.5cm}L{2.1cm}L{4.6cm}@{}}
\toprule
\textbf{App} & \textbf{Seller} & \textbf{Rating (n)} &
\textbf{Price} & \textbf{Notes} \\
\midrule
\endhead
Regimen: Peptide Tracker &
Awaken Labs LLC &
4.8 (338) &
Free; \$4.99/mo, \$39.99/yr &
Health \& Fitness. Updated within the week of retrieval (v2.1.2).
200+ compounds, half-life plotting, labs, 50+ symptom markers,
Apple Health. Health data linked to user ID. \\[0.4em]

DoneDose: Injection Tracker &
Benjamin Fox &
4.0 (5) &
Free; \$4.99/mo, \$34.99/yr, \$4.99 lifetime &
Health \& Fitness. \good{Latest version 1.1.2 dated ``Jun 10'' with no
year shown}---Apple omits the year for the current year, and the app ID
is a 2025-series listing, so this reads as June 2026 and the app is
actively maintained. GLP-1-first
framing; TRT/HRT listed among many. Local storage, no cloud sync. \\[0.4em]

Hormone Helper &
John Patterson &
5.0 (7) &
\$4.99 paid &
\bad{Lifestyle} category, not Health. Last updated September 2024.
Explicitly trans and non-binary framing; educational video content. \\[0.4em]

Transcapsule &
Transcapsule &
--- &
Free w/ premium &
\bad{Abandoned since 2022} per independent
review~\cite{transsolidarity}. Crowdfunded 2020. Promised anonymized
research data sharing. \\[0.4em]

Solace / Euphoria &
Euphoria.LGBT &
--- &
Free &
\bad{Last updated 2022}~\cite{transsolidarity}. Venture-backed.
See \S\ref{sec:euphoria}. \\[0.4em]

MyTransHealth &
--- &
--- &
--- &
\bad{Dead. Last update 2019; site non-functional}~\cite{transsolidarity}. \\[0.4em]

Trans Memo &
--- &
--- &
--- &
\good{Actively maintained} as of April 2026~\cite{transsolidarity}.
HRT reminder and injection tracking. Lightweight. \\
\bottomrule
\end{longtable}

Three observations follow directly from the table and are worth stating
separately because they are easy to read past.

\textbf{The ratings distribution is extreme.} Regimen has 338 ratings.
Every trans-specific product in the census has single digits or no
retrievable count. This is not evidence that trans users do not want
software; it is evidence that they have not been offered a maintained
product to rate.

\textbf{Hormone Helper is in the Lifestyle category.} Its listing is
categorized under Lifestyle rather than Health \& Fitness or
Medical~\cite{hormonehelper}. Whether that was a positioning choice or
an attempt to avoid health-app review scrutiny, the effect is that the
most explicitly trans-framed maintained product in the census is not
present in the category where users search for health tools. It is also,
incidentally, the categorization that would exempt it from Apple's 2027
medical-device declaration requirement (\S\ref{sec:deviceclaim}).

\textbf{DoneDose stores data locally with no cloud sync.} This is the
strongest privacy posture in the census and it belongs to a
GLP-1-oriented product with five ratings. The feature that the
trans market would most reward is currently shipping in the product
least aimed at it.

\subsection{The abandonment pattern}
\label{sec:graveyard}

The trans-specific segment of this category is not a competitive field.
It is a graveyard, and it is documented as one by the community itself.
A 2026 community-authored guide to applications for trans people
characterizes Transcapsule as ``abandoned'' since 2022 and notes that
its website auto-downloads software; MyTransHealth as dead with a
non-functional site since 2019; and Solace as last updated in 2022 with
outdated Android compatibility~\cite{transsolidarity}. Of the sixteen
applications that guide reviews, the ones it recommends without
reservation for trans users are, overwhelmingly, not trans-specific
products at all---Reddit, Finch, Google Maps' location sharing.

This is the central competitive finding of this study, and it inverts
the usual strategic conclusion. The instinct on discovering an empty
segment is to celebrate the opening. The correct reading is more
uncomfortable: the segment is empty because entrants keep dying in it,
and every prospective user of xT has watched that happen at least once.

The practical consequence is that the strongest competitor xT faces is
not a feature set. It is a prior---a learned, evidence-backed
expectation that a trans health app will be built by someone
well-meaning, will be charming for eighteen months, and will then stop
receiving updates while continuing to hold its users' hormone histories
on a server nobody is watching. No feature answers that prior. What
answers it is a name, an entity, and a publication record that together
look like something built to persist.

That reframing has direct nomenclature consequences, developed in
\S\ref{sec:institutional}: it argues for a name with institutional
weight over a name with warmth, because the thing being communicated is
durability rather than empathy.

\subsection{Pharmacokinetic modeling is not a moat}
\label{sec:notamoat}

The project's existing business plan states that xT is ``the only one
(as far as this research found) that's publishing its own PK research
memos and running empirical persona validation against the model,'' and
that ``nobody else in this space appears to be citing population-PK
papers in their source comments.''

The first clause is defensible. The second is not, and the broader
implication---that competitors do not do pharmacokinetics---is false.

Regimen publishes a public Half-Life Visualizer that, by its own
methodology statements, uses ``the Bateman equation, a one-compartment
pharmacokinetic model,'' applies superposition to layer multiple doses
and show steady-state accumulation, and cites FDA NDA labels, a New
England Journal of Medicine trial for retatrutide, and published
literature for testosterone~\cite{regimenhalflife}. It carries
limitation language that would not be out of place in xT's own
documentation: ``Population-average pharmacokinetic model. Individual
variation is significant''; ``real blood levels vary based on individual
metabolism, body composition, injection site, formulation''; and an
acknowledgement that the tool offers ``a useful approximation'' rather
than clinical precision.

That is, to a first approximation, the same model class xT implements,
with the same structural caveats, from a competitor with 338 ratings
and a weekly release cadence.

This correction matters because a positioning built on a false
differentiator fails on contact with a comparison. If the pitch is ``we
model pharmacokinetics,'' a reader who has seen Regimen's visualizer
has already discounted it. The differentiators that survive the
comparison are narrower and, fortunately, real:

\begin{enumerate}[leftmargin=1.5em,itemsep=0.3em]
\item \textbf{Derivation rather than assertion.} Regimen states that it
      uses a Bateman model. xT publishes the closed-form derivation of
      its scaling constant from a named population-PK source, and a
      conservation-of-mass identity for the esters where no primary
      source exists. That is a different epistemic object.
\item \textbf{Published out-of-sample validation.} A validation dossier
      and a persona study reporting error bands against published
      anchors is not something any competitor in this census offers.
\item \textbf{Disclosed defects.} The current revision of xT's PK paper
      reports three defects found by re-deriving the model from its
      sources---including a baseline keyed to the wrong gonads, which
      over-states a transmasculine user's endogenous testosterone by
      more than an order of magnitude at treatment onset---and states
      their magnitudes rather than quietly repairing them. No
      competitor does this. In a market whose defining problem is
      earned distrust (\S\ref{sec:graveyard}), a published defect list
      is worth more than a feature.
\item \textbf{Trans-specific parameterization.} Regimen's compound
      library is organized around peptides, GLP-1 agonists, and TRT.
      Estradiol valerate and cypionate, anti-androgen interactions, and
      transfeminine dosing are, at best, long-tail entries in a product
      whose center of gravity is elsewhere.
\end{enumerate}

The honest one-line claim is therefore not ``we have a PK model.'' It
is closer to: \emph{the only hormone estimator that shows its work, and
the only one that tells you where it is wrong.}

\subsection{What the competitive set implies for the name}

Two conventions are visible in the census and both constrain the brand.

First, the App Store naming convention in this category is
\emph{brand plus descriptor}: ``Regimen: Peptide Tracker,''
``DoneDose: Injection Tracker.'' DoneDose in fact appears under at
least two descriptor variants against a single app
ID~\cite{donedoselisting}, which is ordinary keyword experimentation.
Apple allows 30 characters of app name plus a 30-character
subtitle~\cite{mobileaction}. This is strategically useful: the
descriptor half can carry the search keywords, which frees the brand
half to be abstract. A brand does not have to say ``hormone'' for the
listing to rank for ``hormone tracker.''

Second, and consequently, the argument that an abstract name is
``expensive to market'' because it has to be taught---an argument made
in the project's own monetization document against the byTime
candidate---is weaker than it appears in a store that gives you a
subtitle. The real cost of an abstract name is not comprehension. It is
brand-term search competition, which is a function of collisions, not
of abstractness. That distinction drives the screening in
\S\ref{sec:candidates}.

% ======================================================================
\section{The social, cultural, and political environment}
\label{sec:sociopolitical}

\subsection{Why a naming study contains a legal section}

In most categories, the political environment is a marketing context.
Here it is a threat model. The argument of this section is that the
``discreet'' half of the brand brief is downstream of documented
enforcement mechanisms, and that treating it as an aesthetic preference
would be a category error.

\subsection{Legal climate as of August 2026}
\label{sec:legal}

Two facts establish the direction of travel.

Restrictions have moved from minors to adults, and a federal appellate
court has now permitted it. On 11 March 2026 the United States Court of
Appeals for the Fourth Circuit upheld West Virginia's exclusion of
gender-affirming surgeries for adults from Medicaid coverage, reasoning
that because ``the law applies to specific procedures, not specific
individuals, it does not illegally discriminate against transgender
people''~\cite{advocate4th}. Reporting characterized it as the first
federal appellate enforcement of restrictions on adult
care~\cite{advocate4th}. The reasoning---that a procedure-targeted rule
is not a person-targeted rule---is portable to any restriction framed
around a class of treatment, and hormone therapy is such a class.

State-level data reporting is being proposed. Tennessee House Bill 754
would require clinics to submit patient data to the state including
age, sex, medications, procedures, provider identity, and mental health
diagnoses~\cite{transvitae}. The re-identification objection is not
speculative and has been stated plainly by community outlets: ``If a
report shows that in a small county, one 16-year-old began hormone
therapy in March through a specific provider, it doesn't take a hacker
to figure out who that might be''~\cite{transvitae}.

Coverage protections remain substantially intact in much of the
country---the Movement Advancement Project's healthcare map, updated
25 August 2026, records that 54\% of transgender people live in states
where insurers are prohibited from excluding medically necessary
transgender care, against 2\% of LGBTQ people ages 13+ living in states
that explicitly permit insurers to refuse it~\cite{map}. The picture is
therefore not uniform collapse. It is sharp and worsening divergence,
which for a national consumer product means the brand must be safe in
the worst jurisdiction it ships to, not the median one.

\subsection{Surveillance and health data}
\label{sec:surveillance}

The Project On Government Oversight's investigation into enforcement
technology enumerates the mechanisms: geofence warrants and cellular
location data, data-broker purchases made without warrants, health data
collection and sharing between providers, prescription records, and
driver's license records---noting in one instance a request for more
than 16,000 Texas driver's license gender-change
records~\cite{pogo}. On applications specifically, it warns that
``location data from apps that track our physical whereabouts,
communications through popular apps like Facebook and Instagram, health
data from apps that monitor and track our physical well-being, and
browser search histories paint a detailed picture of our inner
lives''~\cite{pogo}. That investigation is dated June 2023 and predates
the 2026 developments in \S\ref{sec:legal}; it should be read as a floor.

The project's own business plan correctly establishes that HIPAA does
not bind a direct-to-consumer application. The positioning consequence
is the one the plan already draws---trust must be earned by disclosure
rather than by claiming a regulatory badge---but there is a second
consequence the plan does not draw, and it is a naming consequence:

\begin{quote}
\emph{Every artifact bearing the product's name is a potential
disclosure. The App Store purchase record, the recurring line item on a
bank or card statement, the notification banner, the home screen icon
label, the receipt email, and the device backup filename are all
strings the user does not control and cannot redact.}
\end{quote}

A name that encodes ``hormone,'' ``trans,'' ``T,'' or ``HRT'' converts
each of those artifacts into an outing surface. This is the
justification for the discretion test in \S\ref{sec:discretion}, and it
is the single strongest argument against retaining \name{xT} as a
consumer-facing name---not that xT is ugly, but that its one piece of
semantic content is the piece that should not be there.

\subsection{The DIY population}
\label{sec:diy}

A substantial fraction of the addressable population is not
clinician-supervised. Prevalence estimates for do-it-yourself hormone
therapy vary enormously by region and population: a 2022 review reports
figures from 11\% of transgender people in Ontario at the low end to
79\% of trans women in Rio de Janeiro at the high end, with 31\% in
London, 49\% among trans women in San Francisco, and 63\% among trans
people of color in San Francisco~\cite{diywiki}. Reported drivers
include difficulty finding knowledgeable and non-stigmatizing
clinicians, cost, exclusion, discomfort in clinical relationships, and
a desire for faster transition; prevalence is generally higher among
trans women, plausibly because testosterone's controlled-substance
status makes transmasculine DIY harder~\cite{diywiki}.

For xT this cuts two ways, and the tension is worth stating rather than
resolving.

\textbf{It is where the product is most valuable.} A user without
regular bloodwork is exactly the user for whom a superposition model
with a stated error band is the best available estimate of where their
levels are. The PK paper frames the clinical gap as the interval
between blood draws; for this population the interval is unbounded.

\textbf{It is where the product is most dangerous.} The same user has
no clinician to catch the model when it is wrong, and xT's own current
revision documents a defect that over-states endogenous testosterone by
more than an order of magnitude at treatment onset for transmasculine
users. A brand that positions as authoritative to an unsupervised
population is making an implicit promise the model cannot yet keep.

The resolution is not to exclude these users but to let the brand carry
the epistemic status honestly---which favors a name meaning
\emph{estimate}, \emph{model}, or \emph{instrument} over one meaning
\emph{truth}, \emph{health}, or \emph{care}. Names in the latter
family---\name{Verity}, \name{Attest}, \name{TrueLevel} and their
relatives---were generated and discarded on this ground.

\subsection{The Euphoria case}
\label{sec:euphoria}

The most instructive failure in this category is not a technical one.
Euphoria.LGBT (later Solace) launched a suite of trans-focused
applications with venture backing and prominent endorsement, and drew
sustained community criticism. A 2022 critique published by the
Massachusetts Transgender Political Coalition objected on grounds
worth listing individually, because each is a naming and positioning
lesson rather than a product one~\cite{mtpc}:

\begin{enumerate}[leftmargin=1.5em,itemsep=0.25em]
\item \textbf{The tagline.} ``Transition as a service'' was read as
      commodifying transition---treating it as a purchasable product
      rather than a personal process.
\item \textbf{Linearity.} A checkbox-based, staged model of transition
      contradicted lived experience; the author writes that ``no trans*
      person I know has had or wanted a linear transition.''
\item \textbf{Unsourced authority.} Content was published with ``no
      citations or author credentials.''
\item \textbf{Exclusion.} No material specifically by or for trans
      people of color, disabled trans people, or low-income trans
      people.
\item \textbf{Extraction.} Fears of ``data mining and invasions of
      LGBTQ+ privacy,'' and objection to the ethics of venture capital
      in this space specifically.
\item \textbf{Redundancy.} ``All the information they share is already
      being circulated via online and in-person LGBTQ+ communities.''
\end{enumerate}

Three of these bear directly on xT's brand.

\emph{On the name itself.} ``Euphoria'' names the emotional payoff of
transition and sells it back. It is the archetype of the warm,
community-legible name, and it is precisely the register that produced
the commodification objection. This is direct evidence for the
founder's instinct in rejecting the community-native branding axis: in
this category, emotional naming reads as extractive when it comes from
a commercial entity, and the warmth that is supposed to build trust
does the opposite.

\emph{On sourcing.} The ``no citations or author credentials''
objection is the one xT is structurally best positioned to answer, and
it should be answered in the brand rather than in a footnote. A product
that publishes its derivations is immune to the specific criticism that
killed the best-funded entrant in its category.

\emph{On linearity.} The critique of staged, checkbox transition models
is a live risk for xT specifically, because the application ships a
Stage System that binds named clinical scales---Tanner staging among
them---to timeline copy. That feature is defensible as clinical
information and indefensible as a prescription of how transition
should proceed, and the difference is almost entirely a matter of
naming and copy register. It is out of scope here but flagged: the
brand should not imply a finish line.

% ======================================================================
\section{Category positioning}
\label{sec:positioning}

\subsection{The brief as a coordinate, not a checklist}
\label{sec:brief}

The two selected axes---clinical credibility and discretion---define a
plane. Placing the category's actual inhabitants on it is more
informative than any of them individually.

\begin{longtable}{@{}L{3.4cm}L{3.5cm}L{3.5cm}L{4.6cm}@{}}
\toprule
& \textbf{Low discretion} & \textbf{High discretion} & \\
\midrule
\endhead
\textbf{High clinical credibility} &
\emph{Balance} (menopause; clinician-fronted, but unmistakably a
menopause app) &
\bad{Largely vacant.} Regimen approaches it but is coded to male
optimization by its descriptor and compound library. &
The target region. \\[0.6em]

\textbf{Low clinical credibility} &
\emph{Euphoria/Solace}, \emph{Transcapsule}, \emph{Hormone Helper} ---
names that disclose the user &
\emph{Trans Memo}, generic medication reminders &
Where the trans-specific segment actually sits. \\
\bottomrule
\end{longtable}

Two structural observations follow.

\textbf{The target region is vacant for a reason that is not
inaccessibility.} It is vacant because the two axes pull naming in
opposite directions along the dimension of \emph{specificity}. Clinical
credibility is conventionally signalled by naming the domain---
``Hormone Helper'' is legible precisely because it says hormone.
Discretion is achieved by not naming the domain. A name cannot resolve
this by compromise; a half-disclosed name is simply a disclosed name.
It resolves only by relocating the credibility signal off the noun and
onto the \emph{register} of the word---which is what an obscure
technical term does, and what a warm word cannot.

\textbf{This is the whole argument for a technical-register name.} A
term drawn from measurement, mathematics, or instrumentation carries
clinical credibility through its diction while carrying zero
disclosure through its denotation. A clinician reads competence; a
bystander reads nothing. No other naming family satisfies both axes
simultaneously.

\subsection{Institutional weight over warmth}
\label{sec:institutional}

Section \ref{sec:graveyard} established that the dominant competitor is
a learned expectation of abandonment. Brands answer that expectation
with signals of permanence, and the vocabulary of permanence is
institutional: the name of a laboratory, an instrument, an observatory,
a standard. It is not the vocabulary of companionship.

This aligns with, and independently supports, the founder's rejection
of the community-native axis. It also reframes what ``clinically
credible'' is buying. It is not primarily buying provider adoption or
grant eligibility, useful as those are. It is buying the impression
that this is a thing that will still exist in four years---which is the
objection actually forming in a prospective user's mind, whether or not
they articulate it.

\subsection{App Store mechanics}
\label{sec:asomechanics}

The constraints are fixed and favorable to an abstract brand.

\begin{longtable}{@{}L{4.6cm}L{10.6cm}@{}}
\toprule
\textbf{Field} & \textbf{Constraint and consequence} \\
\midrule
\endhead
App name & 30 characters; roughly the first 26 display in iPhone search
results~\cite{mobileaction}. Holds ``\emph{Brand}: \emph{Descriptor}''
comfortably. \\[0.35em]
Subtitle & 30 characters, displayed beneath the name~\cite{mobileaction}.
A second keyword surface. \\[0.35em]
Keywords field & 100 characters, not user-visible, indexed for
search~\cite{mobileaction}. Carries ``hrt, testosterone, estradiol,
injection, trans'' without any of those words appearing on the user's
device. \\[0.35em]
Category & Health \& Fitness is the correct category and triggers the
declaration in \S\ref{sec:deviceclaim}. Hormone Helper's Lifestyle
placement avoids that at the cost of search
presence~\cite{hormonehelper}. \\
\bottomrule
\end{longtable}

The keywords field deserves emphasis because it resolves the apparent
tension between discretion and discoverability. Search terms live in a
field the user never sees and that never appears on their device. A
maximally discreet brand name loses nothing in ranking for ``hormone
tracker'' provided the descriptor and keyword fields are used. The
discretion tax on discoverability is close to zero; the discretion tax
falls entirely on brand-term recall, which is a marketing cost, not a
search cost.

\subsection{The claim boundary: Apple's medical device declaration}
\label{sec:deviceclaim}

This subsection sits outside the requested scope and is included
because it constrains what the brand may assert.

As of 26 March 2026, Apple requires that an app be assigned a
regulated-medical-device status in App Store Connect if its primary or
secondary category is Health \& Fitness or Medical, or if it is marked
as containing frequent references to medical or treatment information
in the age-rating questionnaire~\cite{appledevice}. Apple defines
regulated medical device apps as ``those that function on their own or
as part of a system for a range of medical purposes, including
diagnosis, prevention, \emph{monitoring}, and treatment of diseases and
physiological conditions'' (emphasis added). New apps must declare from
that date; existing apps in the EEA, UK, and US must declare by early
2027, after which ``you'll no longer be able to submit app
updates''~\cite{appledevice}.

xT will be a new app in Health \& Fitness and will therefore declare at
submission. The declaration itself is routine---an app that is not a
regulated device selects No. The positioning consequence is the
constraint that comes with it: the marketing register and the declared
status must agree. A brand and tagline built on \emph{prediction of
physiological state} sit closer to ``monitoring'' than a brand built on
\emph{logging and estimation} does. This argues against names and
taglines in the diagnostic register---\name{Assay}, \name{Verity}, and
anything reading as a measurement of the user rather than a model of a
schedule---independently of their trademark profiles.

The safe formulation is the one the project has already written into
its own disclaimer: xT estimates hormone levels from a pharmacokinetic
model; it does not replace bloodwork or medical advice. The brand
should be consistent with that sentence, which means the name should
denote a \emph{model} or a \emph{curve}, not a \emph{measurement} or a
\emph{verdict}.

\subsection{The position to claim}

Synthesizing \S\ref{sec:notamoat}, \S\ref{sec:graveyard},
\S\ref{sec:brief}, and \S\ref{sec:deviceclaim}:

\begin{quote}
\emph{A maintained, technically serious instrument for people managing
hormone therapy---one that publishes how it works and where it is
wrong, names nothing about its user, and is built to still be here in
four years.}
\end{quote}

Every clause is load-bearing. \emph{Maintained} answers the graveyard.
\emph{Publishes how it works and where it is wrong} is the only
differentiator that survives comparison with Regimen and the only one
that answers the Euphoria critique. \emph{Names nothing about its user}
is the safety requirement. \emph{Instrument} rather than \emph{monitor}
keeps the claim inside the declaration boundary.

Notably, the phrase ``for trans people'' does not appear---not because
the product is not for trans people, which it emphatically is and
should say everywhere except on the user's device, but because the
brand layer is the one layer that has to travel onto a lock screen in a
hostile state. Trans-specificity belongs in the App Store description,
the landing page, the community posts, and the grant applications. It
does not belong in the string.

% ======================================================================
\section{Rebrand}
\label{sec:rebrand}

\subsection{The discretion test}
\label{sec:discretion}

``Discreet'' is not a feeling a name gives; it is a property that can be
tested. Four tests operationalize it, in rough order of how often each
situation occurs.

\begin{longtable}{@{}L{3.2cm}L{11.9cm}@{}}
\toprule
\textbf{Test} & \textbf{Question} \\
\midrule
\endhead
Lock screen &
A notification arrives while the phone is face-up on a desk, or is read
over a shoulder on transit. Does the app name, plus a plausible
notification body, disclose that the user is on hormone therapy?
Note that the notification \emph{copy} is as much a disclosure surface
as the name---``Time for your injection'' passes; ``Estradiol dose
due'' does not. \\[0.4em]

Financial record &
A recurring subscription line item appears on a bank or card statement,
a shared family plan, or an expense report. Does the descriptor
disclose? This surface is not fully controllable: the App Store
receipt shows the app name as listed. \\[0.4em]

Over the shoulder &
The home screen icon and its label are visible to a partner, parent,
employer, roommate, or border officer. Does the icon disclose---by
wordmark, by symbol, or \emph{by palette}? \\[0.4em]

Compelled disclosure &
A device is examined under legal compulsion or coercion. Nothing
protects the user here except that the name and icon are not
self-incriminating on their face, and that the vendor holds as little
as possible. This test cannot be passed by naming alone; it is listed
because it sets the ceiling on what naming can accomplish and argues
for the data-minimization posture separately. \\
\bottomrule
\end{longtable}

The third test has an implication the founder's own design language
currently fails, and it is worth stating plainly because it is
otherwise easy to miss in a naming document. The existing visual
identity---accent \code{\#E8709A} rose against \code{\#7E74D8}
lavender---is, in 2026, a legible trans signifier. A rose-and-lavender
glass icon on a home screen discloses roughly as much as the word
``hormone'' would. Discretion is a property of the mark, not only of
the word, and \S\ref{sec:wordmark} treats it accordingly.

\subsection{Screening method}

Each candidate was screened on five criteria:

\begin{description}[leftmargin=1.4em,style=nextline,itemsep=0.25em]
\item[Discretion] Against the four tests above.
\item[Clinical register] Whether the word's diction reads as technical
      rather than consumer, without asserting a diagnostic claim
      (\S\ref{sec:deviceclaim}).
\item[App Store collision] Whether a storefront search for the bare
      word returns competing results that would bury the brand.
\item[Trademark screen] United States federal records via Justia,
      filtered to classes 9, 42, and 44, with class 5
      (pharmaceuticals) noted where relevant. \emph{Indicative only;
      see \S\ref{sec:cannot}.}
\item[Sayability] Whether the word survives being said aloud to a
      clinician, read by someone who has never seen it, and typed from
      memory.
\end{description}

\subsection{Screened candidates}
\label{sec:candidates}

\subsubsection*{1. \name{Isocline} --- recommended}

\emph{An isocline is a curve along which a quantity holds constant---in
a differential system, the locus of points sharing the same slope.}

\begin{longtable}{@{}L{3.0cm}L{12.1cm}@{}}
\toprule
\endhead
Rationale &
It names the mathematics of the product rather than its subject
matter. The application's entire function is the shape of a curve over
time, and an isocline is the object one draws when asking where a
system is heading rather than where it is. It denotes a \emph{model},
not a \emph{measurement}, which keeps it on the correct side of the
claim boundary in \S\ref{sec:deviceclaim}. \\[0.35em]
Discretion & \good{Passes all four.} Discloses nothing; reads as
geology, mathematics, or a consulting firm. \\[0.35em]
Register & \good{Strong.} Unambiguously technical diction; carries
credibility without a health noun. Works equally as a product name and
as the name of a research entity (``Isocline Research''), which matters
for the eventual nonprofit arm. \\[0.35em]
App Store & \good{Clear.} A storefront search surfaced no competing
application. \\[0.35em]
Trademark & \good{Clear in the relevant classes.} Four marks located:
Turin Golf (class 28, 2000), Huron Graffenstaden (class 7, 1991),
Centre d'Etudes du Fraisage (class 7, 1982), and Isocline Ventures LLC
(class 36, financial services, 2014). \good{Nothing in classes 9, 42,
44, or 5.} The three industrial and sporting marks are old enough that
their live status should be confirmed but are in unrelated classes
regardless. \\[0.35em]
Domain & \mixed{Assume \code{isocline.com} is held} by Isocline
Ventures LLC. \code{isocline.app} and \code{isocline.health} are the
realistic targets. \bad{Unverified---check at a registrar.} \\[0.35em]
Weaknesses & Four syllables and a spelling that has to be heard twice.
Semantically precise but not self-explanatory---the meaning has to be
told, once, in the About screen. That is a real cost and it is the main
argument for \name{Ephemeris} or a shorter name instead. \\
\bottomrule
\end{longtable}

\subsubsection*{2. \name{Ephemeris} --- principal alternate}

\emph{An ephemeris is a table giving the computed future positions of a
body at regular intervals---the oldest form of the thing xT is.}

\begin{longtable}{@{}L{3.0cm}L{12.1cm}@{}}
\toprule
\endhead
Rationale &
The semantic fit is the best of any candidate considered. An ephemeris
is precisely a precomputed trajectory: not a measurement of where
something is, but a principled prediction of where it will be, tabulated
against time, produced from a model and periodically corrected against
observation. That is a one-word description of a superposition PK
engine with blood-test calibration. It also carries genuine
institutional gravity---ephemerides are published by observatories, and
have been for centuries, which is exactly the durability signal
\S\ref{sec:institutional} calls for. \\[0.35em]
Discretion & \good{Passes all four.} \\[0.35em]
Register & \good{Strong}, with one caveat: astrological ephemerides
are the more common consumer encounter with the word, and that
association cuts against clinical credibility for some readers while
being neutral-to-charming for others. \\[0.35em]
App Store & \mixed{Contested but not in health.} ``Ephemeris: Moon \&
Sun Tracker,'' ``iPhemeris Astrology,'' and ``The Photographer's
Ephemeris'' occupy the term. None is a health application, so category
browsing is unaffected, but a bare brand-term search is crowded. \\[0.35em]
Trademark & \mixed{Moderate.} Approximately 30 records. Two in class 9:
Moonshadow Mobile (2014) and The Photographer's Ephemeris / Crookneck
Consulting (2015). Others in classes 33, 25, 14, 3. \good{None in class
42 or 44.} The class 9 marks are the concern and warrant an attorney's
look. \\[0.35em]
Domain & \bad{Unverified.} Assume the \code{.com} is gone. \\[0.35em]
Weaknesses & Long, hard to spell from hearing, and the astrology
association is a genuine liability against the ``clinically credible''
half of the brief---the one place where this candidate is weaker than
\name{Isocline} despite the better semantics. \\
\bottomrule
\end{longtable}

\subsubsection*{3. \name{Bateman} --- best clearance, disqualifying association}

\emph{Harry Bateman, whose equation the application implements.}

\begin{longtable}{@{}L{3.0cm}L{12.1cm}@{}}
\toprule
\endhead
Rationale &
Eponymous names carry more institutional weight than any other family:
they read as a laboratory, a standard, or a firm, and they are the
strongest available answer to the abandonment prior in
\S\ref{sec:graveyard}. It is also the most honest name on this list ---
the product is a Bateman-function engine. \\[0.35em]
Discretion & \good{Passes all four} completely; a surname discloses
nothing. \\[0.35em]
Trademark & \good{The cleanest screen of any candidate.} Only two
records with the exact mark: Compass Group Holdings (class 43,
catering, 1994) and Modern Fabrication Inc.\ (material handling, 2019).
\good{Nothing in classes 5, 9, 42, or 44.} \\[0.35em]
Weaknesses & \bad{Patrick Bateman.} The protagonist of \emph{American
Psycho} is, for the exact demographic this product serves---young,
online, extremely fluent in reference---the first and possibly only
association the name will produce. A trans health application named
after a fictional misogynist serial killer is a headline, and the
community documented in \S\ref{sec:euphoria} is more than capable of
writing it. This is not a reason to hesitate; it is a reason to
decline. Listed because its clearance profile is genuinely the best
found and the reasoning should be on the record. \\
\bottomrule
\end{longtable}

\subsubsection*{4. \name{Titer} --- clinical, with a pronunciation problem}

\emph{A titer is the measured concentration of a substance in a
solution.}

\begin{longtable}{@{}L{3.0cm}L{12.1cm}@{}}
\toprule
\endhead
Rationale & Short, real, clinical, ungendered, and directly on-subject
without naming a hormone. \\[0.35em]
Discretion & \good{Passes all four} in writing. \\[0.35em]
Register & \mixed{Strong but on the wrong side of the claim boundary.}
A titer is a \emph{measurement}. xT does not measure; it estimates. A
name asserting measurement sits closer to the ``monitoring'' language
in Apple's device definition (\S\ref{sec:deviceclaim}) and over-claims
against the product's own disclaimer. \\[0.35em]
App Store & \good{Clear.} No competing application surfaced. \\[0.35em]
Trademark & \mixed{Crowded in the adjacent space.} Roughly 90 records,
59 in class 9---but predominantly laboratory-apparatus compounds
(TITERMAX/Promega, TITER ON CHIP/InDevR, TITERSEAL, TITERTUBE) rather
than the bare mark. Dense enough in scientific class 9 to make
clearance uncertain. \\[0.35em]
Weaknesses & \bad{Said aloud, ``titer'' is ambiguous between
\emph{TIE-ter} and a first syllable that reads as ``tit.''} For a
trans health application---and specifically for a transfeminine user
saying the app's name to a clinician, a pharmacist, or a
partner---this is not a minor infelicity. It is a name that can
embarrass its user in exactly the interaction the product exists to
support. Recommended against on this ground alone. \\
\bottomrule
\end{longtable}

\subsubsection*{5. \name{Interval} --- the plainest defensible option}

\emph{The dosing interval; also the gap between blood draws that the
product exists to fill.}

\begin{longtable}{@{}L{3.0cm}L{12.1cm}@{}}
\toprule
\endhead
Rationale & The naming rationale writes itself out of the project's own
pharmacokinetic paper, which frames the clinical gap as the period
\emph{between} draws where patients have no visibility. \name{Interval}
names that gap. It is also the easiest candidate to spell, say, and
remember. \\[0.35em]
Discretion & \good{Passes all four.} Reads as a timer or a fitness app,
which is close to perfect camouflage. \\[0.35em]
Register & \mixed{Neutral.} Plain rather than technical; carries less
credibility than \name{Isocline} or \name{Ephemeris}. \\[0.35em]
App Store & \bad{Poor.} ``Interval'' in a storefront search returns
interval timers and HIIT applications in volume. Brand-term search
would be unwinnable. \\[0.35em]
Trademark & \mixed{No bare-mark conflict located in classes 9, 42, or
44}---the field is dominated by interval-training marks in class 41 and
by Interval International in travel. The generic-word problem is
larger than the conflict problem: a bare common noun is weak as a mark
and hard to enforce. \\[0.35em]
Weaknesses & Wins on plainness, loses on both halves of the brief ---
insufficiently distinctive to own, insufficiently technical to signal
seriousness. \\
\bottomrule
\end{longtable}

\subsubsection*{6. \name{Plateau} --- rejected on semantics, not clearance}

\emph{The steady-state plateau a repeated-dose regimen converges to.}

Pharmacokinetically apt: reaching plateau is what a titrated regimen is
for. Clearance is mediocre---284 records, including Plateau Systems in
class 42 (software, 2001) and Menicon in class 9---but survivable.

It fails on meaning. In the vocabulary of transition, ``plateau'' does
not mean \emph{therapeutic steady state}; it means \emph{stopped
changing}, and it is a word users encounter in anxious forum posts
about whether their development has stalled. Naming the application
after the fear it is meant to allay is a fault no visual identity
repairs. Retained in this list rather than in \S\ref{sec:rejected}
because the clearance work was done and the semantic failure is the
interesting part.

\subsubsection*{7. \name{xT} --- the incumbent}
\label{sec:incumbent}

The founder did not select ``keep xT'' as a brand posture. It is
benchmarked here anyway, because a rebrand recommendation that has not
beaten the incumbent on the record is not a recommendation. The case
for retention is real and is stated at full strength.

\textbf{In favor.} It is two characters, which is a genuine asset---it
fits any lockup, any icon, any 30-character App Store name with 27
characters left for the descriptor. It is already the name in the
codebase, the repository, the running heads of two research papers, and
the author's own usage; renaming has a real cost measured in
documentation, continuity, and the small credibility that accrues to a
name that has been used consistently. It is ungendered in the sense
that mattered originally: it does not commit the product to estrogen or
to testosterone. It has no plausible trademark conflict, because
two-letter marks in unrelated classes rarely collide in practice. And
short alphanumeric names age well in software.

\textbf{Against.} The decisive objection is not aesthetic. It is that
xT's only semantic content is the wrong content. The ``T'' reads as
testosterone---the founder's own business plan says exactly this, that
``\name{xT} at least gestured at testosterone/hormones''---and that
gesture is the single piece of information the discretion analysis in
\S\ref{sec:surveillance} says must not be on the device. A name that
is otherwise opaque, but whose one legible element is the hormone,
fails the lock-screen and over-the-shoulder tests specifically for the
users in the worst jurisdictions.

Three further objections are weaker but real. It is unsearchable: a
two-character string cannot rank or be recalled, and cannot be spoken
to a friend as a recommendation, which matters disproportionately in a
market whose only real acquisition channel is peer recommendation. It
reads as unfinished, because it is---a codename signals a project, not
an institution, which is the opposite of the durability signal
\S\ref{sec:institutional} identifies as the category's central
requirement. And it has no room for a research arm: ``xT Research''
does not read as an organization that publishes.

\textbf{Verdict.} Retention is defensible for a private beta and
indefensible at public launch. If the name changes at all, it should
change before the App Store listing, the domain, the privacy policy,
and the entity filing---all of which are imminent per the project's own
90-day plan, and all of which are more expensive to change afterward
than the name is to change now.

\subsection{Rejected candidates}
\label{sec:rejected}

Eleven names were generated and discarded. The reasons are more useful
than the names, because each identifies a failure mode that will recur
in any further naming work.

\begin{longtable}{@{}L{2.9cm}L{3.6cm}L{8.5cm}@{}}
\toprule
\textbf{Name} & \textbf{Failure mode} & \textbf{Detail} \\
\midrule
\endhead
\name{Sextant} &
\bad{Discretion, catastrophically} &
Semantically excellent---an instrument for fixing your position when
you cannot see land. The string begins with \code{SEX}. For a product
whose brief is that the name must be safe on a lock screen and a bank
statement, this is disqualifying on sight. Also crowded: 178 records
including Sextant Medical in class 44. \\[0.4em]

\name{Serum} &
Collision & Elegant camouflage---cosmetic to a layperson, clinical to a
clinician. But the App Store holds several SERUM applications, and Xfer
Records' Serum is a widely known audio product almost certainly
registered in class 9. \\[0.4em]

\name{Depot} &
Collision & A depot injection is exactly the route modeled. Home Depot
makes the brand term unsearchable. \\[0.4em]

\name{Almanac} &
Collision & Semantically strong (periodic quantities tabulated across
time). ``Almanac Health'' already exists in the App Store---a direct
collision in the exact category. \\[0.4em]

\name{Assay} &
Trademark + claim &
Clean in the App Store, but a live bare \name{ASSAY} mark exists in
class 44 (medical services), owned by Assay LLC, alongside class 9
marks held by PerkinElmer and others. Independently, ``assay'' asserts
measurement and sits on the wrong side of the claim boundary
(\S\ref{sec:deviceclaim}). \\[0.4em]

\name{Isochron} &
Trademark, in pharma &
Better semantics than \name{Isocline} for a time-domain model. But five
class 5 pharmaceutical registrations exist (Forest Laboratories, Sun
Pharmaceutical), plus class 9 and 42 marks. \name{Isochron} is a drug
brand name. Unusable for a health product. \\[0.4em]

\name{Nadir} &
Valence & The trough of a dosing cycle. 169 records including a class
42 filing in February 2026. Independently, naming the product after the
lowest point is the wrong promise. \\[0.4em]

\name{byTime} &
Collision (previously identified) &
The project's prior rename candidate. Two existing iOS applications
share the name, \code{bytimeapp.com} is taken, and the GitHub
organization is claimed. The project's own monetization document
already reached this conclusion; it is repeated here so the record is
complete. \\[0.4em]

\name{Euphoria}-family &
Register &
Warm, emotional, transition-naming words. \S\ref{sec:euphoria}
documents what happened to the best-funded product that tried this.
Excluded by the founder's brief and independently by the evidence. \\[0.4em]

\name{Verity}, \name{Attest}, \name{TrueLevel} &
Over-claim &
Names asserting truth or accuracy. The product's own published
revision documents three defects including an order-of-magnitude
baseline error. A brand may not promise what the model has disclosed it
cannot deliver. \\[0.4em]

\name{Half-Life} &
Trademark & Valve Corporation. No further analysis required. \\
\bottomrule
\end{longtable}

\subsection{Wordmark and visual direction}
\label{sec:wordmark}

Naming solves three of the four discretion tests. The
over-the-shoulder test is solved by the mark, and the current identity
does not solve it.

\subsubsection*{The palette problem}

The existing design language---\code{\#E8709A} rose and \code{\#7E74D8}
lavender on a near-black \code{\#0A0814}, in glass cards---is
well-executed and is the correct emotional register for the product's
\emph{interior}. As an \emph{icon}, it is a disclosure. Rose-and-lavender
gradients are, in 2026, among the most legible visual signifiers of
trans-adjacent software, and a home screen is a public surface.

The recommendation is a split identity, which is a normal pattern and
costs nothing:

\begin{longtable}{@{}L{3.4cm}L{11.7cm}@{}}
\toprule
\textbf{Surface} & \textbf{Treatment} \\
\midrule
\endhead
Icon and lock screen &
\good{Neutral.} Monochrome or near-monochrome: graphite, slate,
bone, or deep ink. A single geometric mark. No gradient, no rose. This
is the only surface where discretion outranks warmth, and it should be
designed as an instrument's badge rather than as an app icon. \\[0.4em]
In-app &
\good{Keep the existing palette.} Once the phone is unlocked and the
app is open, the user has already chosen to be there. The rose and
lavender belong here, and they are doing real work---the softness is
what makes an honest, sometimes unwelcome data presentation bearable. \\[0.4em]
Landing page and store listing &
Warm, and explicitly trans-first in copy. These surfaces are sought
out, not encountered. \\[0.4em]
Research papers and entity &
\good{Austere.} Serif, black on white, no accent color. The papers are
the credibility asset; they should look like the literature they cite. \\
\bottomrule
\end{longtable}

\subsubsection*{Typography}

For \name{Isocline} specifically: a grotesque with a technical rather
than a friendly cast---something in the register of a scientific
instrument's faceplate. Avoid geometric humanists (the Poppins/Circular
family), which read as consumer wellness and undercut the credibility
the name is buying. Set the wordmark in lowercase or small caps with
moderately open tracking; the word is long enough that tight tracking
makes it unreadable at icon scale. Do not letterspace it into an
acronym.

\subsubsection*{The mark}

The obvious symbol is the curve itself, and it is worth resisting the
obvious form. A literal plotted PK curve is a chart, not a mark, and
becomes mush below 60 pixels. The extractable idea is that an isocline
is a \emph{family of parallel slopes}---a small field of short parallel
strokes at a constant angle, which is both what the mathematics
depicts and a shape that survives to 16 pixels. It reads as abstract
and technical at every size and discloses nothing.

The landing page's existing background---a superposed multi-dose curve
taken from the model's own output---should stay. It is genuinely good,
and it belongs on the surface where the product is explaining itself
rather than on the surface where it is sitting on a stranger's desk.

\subsubsection*{Naming lockup}

Following the category convention (\S\ref{sec:asomechanics}), with the
30-character limit in mind:

\begin{itemize}[leftmargin=1.5em,itemsep=0.2em]
\item App Store name: \code{Isocline: Hormone Tracker} (25 characters).
\item Subtitle: something in the register of \code{Levels, doses, and
      lab results} (30 characters) --- deliberately not naming a
      hormone, since the subtitle appears in the store listing but the
      \emph{name} is what propagates to the device.
\item Keywords field: \code{hrt, testosterone, estradiol,}
      \code{injection, trans, hormone, transition, taper} ---
      invisible to the user, indexed by Apple, and present on no
      surface the user carries.
\item Home screen label: \code{Isocline}.
\end{itemize}

Note the deliberate asymmetry: the store listing is findable and
explicit, the device is silent. That asymmetry is the entire design.

\subsection{Recommendation}

\begin{longtable}{@{}L{1.0cm}L{2.8cm}L{11.3cm}@{}}
\toprule
\textbf{Rank} & \textbf{Name} & \textbf{Basis} \\
\midrule
\endhead
1 & \good{\name{Isocline}} &
Only candidate clear on all five screening criteria. Passes all four
discretion tests; unambiguously technical register; no App Store
collision; no located trademark conflict in classes 5, 9, 42, or 44;
denotes a model rather than a measurement. Costs one sentence of
explanation, once. \\[0.4em]
2 & \name{Ephemeris} &
Better semantics, worse everything else: crowded storefront, two class
9 marks, and an astrology association that works against the
credibility half of the brief. Take it if the meaning matters more than
the clearance. \\[0.4em]
3 & \name{xT} (retain) &
Defensible for private beta. Fails the discretion tests at public
launch for the users who need them most. \\[0.4em]
--- & \name{Titer}, \name{Interval}, \name{Plateau}, \name{Bateman} &
Each fails one criterion decisively---pronunciation, distinctiveness,
valence, and cultural association respectively. Documented rather than
recommended. \\
\bottomrule
\end{longtable}

\noindent
Three actions follow, in this order, and the first two are cheap:

\begin{enumerate}[leftmargin=1.5em,itemsep=0.25em]
\item Verify \code{isocline.app} / \code{isocline.health} at a
      registrar and reserve the App Store name in App Store Connect.
      Both are same-afternoon tasks and both are prerequisites for the
      entity filing.
\item Have a trademark attorney run a real clearance search on
      \name{Isocline} in classes 9 and 44 before anything is printed or
      filed. The screen in this paper is not that search
      (\S\ref{sec:cannot}).
\item Only then: entity name, privacy policy, landing page
      find-and-replace, and App Store listing. The landing page rename
      is, as the project's own monetization document notes, a
      find-and-replace.
\end{enumerate}

% ======================================================================
\section{Limitations}
\label{sec:limitations}

\textbf{The trademark screens are not clearance.} Stated in
\S\ref{sec:cannot} and repeated here because it is the finding most
likely to be over-read. Every ``clear'' in \S\ref{sec:candidates}
means only that a Justia-indexed search of United States federal
records surfaced no conflict in the named classes. State registrations,
common-law marks, foreign filings, and likelihood-of-confusion analysis
are all outside it.

\textbf{Domain and App Store name availability were not established.}
No registrar or App Store Connect query was possible with the tooling
available.

\textbf{App Store figures are a single snapshot.} Rating counts,
prices, and last-updated dates were retrieved in August 2026 from
storefront listings. Some competitor data---particularly for the
abandoned applications---comes from a community-authored secondary
source~\cite{transsolidarity} rather than from listings, because the
listings are gone or were not locatable. Those abandonment claims are
consistent with the other evidence but were not independently verified
listing-by-listing.

\textbf{The positioning map in \S\ref{sec:brief} is a judgment, not a
measurement.} No user research, survey, or interview informed this
paper. Every claim about what trans users want is inferred from
published community writing and from competitor outcomes. That is a
real weakness: the single highest-value next step for this work is
twenty conversations with prospective users, which would test the
central claim of \S\ref{sec:graveyard}---that abandonment fear
dominates feature preference---directly rather than by inference.

\textbf{The medical and clinical axis is absent by scope.} Provider
adoption, WPATH and Endocrine Society positioning, and the
software-as-a-medical-device question were deprioritized. Only the
App Store declaration was carried in (\S\ref{sec:deviceclaim}), and
only because it bounds the brand claim.

\textbf{Two prominent sources are competitor marketing} and are used
for enumeration only (\S\ref{sec:method}).

% ======================================================================
\section{Conclusion}

The question was whether to rename xT. The more useful finding was what
the market turned out to be.

It is not a crowded field of HRT trackers that xT must out-feature. It
is a segment whose trans-specific entrants are, with one exception,
dead---Transcapsule since 2022, Solace since 2022, MyTransHealth since
2019---while the maintained, well-rated products adjacent to it are
built for cisgender men optimizing testosterone and peptides. The
opening is real. So is the reason for it: users in this category have
learned that tools built for them get abandoned, and that prior is a
more formidable competitor than any feature set.

The project's own business plan identified the pharmacokinetic model as
the moat. It is not: Regimen ships a one-compartment Bateman
superposition visualizer with cited sources and stated limitations. The
asset that survives comparison is narrower and better---xT publishes
its derivations, its out-of-sample validation, and, uniquely, its own
defects. In a market defined by earned distrust, a published defect
list outperforms a feature.

The name follows from that. It must signal durability rather than
warmth, because durability is the doubt. It must signal technical
seriousness through diction rather than through a health noun, because
the health noun is a disclosure. It must denote a model rather than a
measurement, because that is what the product is and what its own
disclaimer says. And it must survive appearing on a lock screen, a bank
statement, and a home screen in a state that is, as of March 2026,
permitted to restrict adult care.

\name{Isocline} satisfies all of these and cleared every screen this
study could run. \name{Ephemeris} is the better metaphor and the worse
brand. \name{xT} remains adequate for a private beta and should not
survive the App Store listing---not because it is a poor name, but
because its one legible character names the hormone, and that is the
one thing the string must not do.

The visual identity needs the same split the name does: austere on
every surface the user carries in public, and warm everywhere the user
has chosen to be. The rose and the lavender are not the problem. Their
location is.

% ======================================================================
\begin{thebibliography}{99}

\bibitem{advocate4th}
Advocate (2026). ``U.S.\ court allows state bans on gender-affirming
care for adults in unprecedented ruling.'' Reporting on the U.S.\ Court
of Appeals for the Fourth Circuit decision of 11 March 2026 upholding
West Virginia's Medicaid exclusion for adult gender-affirming surgery.
\url{https://www.advocate.com/news/court-bans-adult-genderaffirming-care}

\bibitem{appledevice}
Apple Inc.\ (2026). ``Update on regulated medical device apps in the
European Economic Area, United Kingdom, and United States.'' Apple
Developer News, 26 March 2026.
\url{https://developer.apple.com/news/?id=nyqbfz1y}

\bibitem{diywiki}
``DIY transgender hormone therapy.'' Wikipedia. Prevalence figures
cited therein to a 2022 review; region-specific estimates range from
11\% (Ontario) to 79\% (Rio de Janeiro).
\url{https://en.wikipedia.org/wiki/DIY_transgender_hormone_therapy}

\bibitem{donedoseblog}
Done Dose (2026). ``Best HRT Tracker Apps for Hormone Replacement
Therapy (2026).'' \bad{Interested source} --- published by DoneDose,
which appears in its own ranking.
\url{https://www.donedose.com/guides/best-hrt-tracker-app}

\bibitem{donedoselisting}
``DoneDose: Injection Tracker'' / ``DoneDose: Shot \& Med Tracker,''
Apple App Store, app ID 6757083101. Seller: Benjamin Fox. Retrieved
August 2026.
\url{https://apps.apple.com/us/app/donedose-injection-tracker/id6757083101}

\bibitem{hormonehelper}
``Hormone Helper,'' Apple App Store, app ID 6517361645. Seller: John
Patterson. Lifestyle category. Retrieved August 2026.
\url{https://apps.apple.com/us/app/hormone-helper/id6517361645}

\bibitem{map}
Movement Advancement Project. ``Healthcare Laws and Policies''
equality map. Last updated 25 August 2026.
\url{https://mapresearch.org/equality-map/healthcare-laws-and-policies}

\bibitem{mobileaction}
MobileAction (2026). ``How to name an app for App Store \& Google
Play: 2026 Guide.'' Source for the 30-character app name limit, the
30-character subtitle, and the 100-character keywords field.
\url{https://www.mobileaction.co/blog/how-to-name-an-app/}

\bibitem{mtpc}
Carnell H.\ (2022). ``There is No Euphoria Here.'' Massachusetts
Transgender Political Coalition, 23 June 2022.
\url{https://www.masstpc.org/there-is-no-euphoria-here/}

\bibitem{pogo}
Project On Government Oversight (2023). ``Policing Gender: How
Surveillance Tech Aids Enforcement of Anti-Trans Laws.'' 28 June 2023.
\url{https://www.pogo.org/investigates/policing-gender-how-surveillance-tech-aids-enforcement-of-anti-trans-laws}

\bibitem{regimenblog}
Regimen (2026). ``Best HRT Tracker Apps 2026 (6 Compared).''
\bad{Interested source} --- published by Regimen, which appears first
in its own ranking.
\url{https://helloregimen.com/blog/best-hrt-tracker-apps-2026}

\bibitem{regimenhalflife}
Regimen. ``Half-Life Visualizer: Testosterone \& Peptide Blood
Levels.'' Methodology statements quoted in \S\ref{sec:notamoat}.
\url{https://helloregimen.com/tools/half-life-visualizer}

\bibitem{regimenlisting}
``Regimen: Peptide Tracker,'' Apple App Store, app ID 6753905449.
Seller: Awaken Labs LLC. Retrieved August 2026.
\url{https://apps.apple.com/us/app/regimen-peptide-tracker/id6753905449}

\bibitem{transsolidarity}
Trans Solidarity Project (2026). ``Essential Apps for Transgender
People in 2026.'' Community-authored review; source for the
abandonment status of Transcapsule, MyTransHealth, and Solace, and for
Trans Memo's continued maintenance as of April 2026.
\url{https://transsolidarityproject.org/transgender-apps-2026/}

\bibitem{transvitae}
TransVitae. ``When Trans Healthcare Data Becomes Government
Surveillance.'' Source for Tennessee House Bill 754 and the
re-identification argument.
\url{https://www.transvitae.com/when-trans-healthcare-data-becomes-government-surveillance/}

\bibitem{justia}
Justia Trademarks. United States federal trademark records, searched
August 2026 for: assay, titer, isocline, isochron, ephemeris, sextant,
plateau, nadir, bateman, interval. Filtered to International Classes 9,
42, and 44, with class 5 noted where relevant.
\url{https://trademarks.justia.com/}

\end{thebibliography}

\end{document}
